


Use this guide to upload grades for your assigned subjects, correct individual records, and download grade reports. Grade results are visible to students only after the Director triggers a publish.

<Info>
  You log in to UniZ with the same endpoint used by admins. Your `role` in the
  login response determines your access level.
</Info>

\subsection{Upload grades via Excel template}

The recommended workflow for submitting an entire batch's grades is to download the pre-filled Excel template, fill in grades, and upload the completed file.


\section{Instruction Step}

  
\section{Instruction Step}

    Send a `GET` request to `/academics/grades/template` with query parameters
    for the batch you are grading. The template comes pre-filled with all
    student IDs for that batch, so you do not need to enter them manually.

    \begin{lstlisting}[language=bash]
    GET https://api.uniz.rguktong.in/api/v1/academics/grades/template?branch=CSE&year=E3&semester=SEM-1&subjectCode=E3-SEM-1-CSE-1
    Authorization: Bearer <your-token>
    
\end{lstlisting}

    | Parameter     | Example             | Description                          |
    | ------------- | ------------------- | ------------------------------------ |
    | `branch`      | `CSE`               | Department code                      |
    | `year`        | `E3`                | Batch year (E1–E4, P1–P2)            |
    | `semester`    | `SEM-1`             | Semester identifier                  |
    | `subjectCode` | `E3-SEM-1-CSE-1`    | Subject code as registered in UniZ   |

    The response is an `.xlsx` file download. Save it locally before filling it in.

    <Tip>
      In the UniZ portal, you can download the template directly from the
      **Add Grades** page without writing the request manually.
    </Tip>

  

  
\section{Instruction Step}

    Open the downloaded template. Each row corresponds to one student. Fill in
    the **Grade** column using the accepted grade values:

    | Value  | Meaning        |
    | ------ | -------------- |
    | `EX`   | Excellent       |
    | `A`    | Grade A         |
    | `B`    | Grade B         |
    | `C`    | Grade C         |
    | `D`    | Grade D         |
    | `F`    | Fail            |
    | `9.5`  | Numeric points  |

    Do not modify the student ID column or add or remove rows. Leave the grade
    cell blank only if a student was absent for evaluation — do not use `N/A`
    or similar placeholders.

    <Warning>
      The file size limit for uploads is **50 MB**. Standard grade sheets are
      well within this limit, but avoid embedding images or charts in the
      spreadsheet.
    </Warning>

  

  
\section{Instruction Step}

    Post the completed file to `/academics/grades/upload` as a multipart form.

    \begin{lstlisting}[language=bash]
    POST https://api.uniz.rguktong.in/api/v1/academics/grades/upload
    Authorization: Bearer <your-token>
    Content-Type: multipart/form-data

    file=<your-filled-template.xlsx>
    
\end{lstlisting}

    **Success response:**

    \begin{lstlisting}[language=json]
    {
      "success": true,
      "processed": 62,
      "successCount": 60,
      "failCount": 2,
      "errors": [
        "Row 14: Invalid grade value 'NA' for student O210045",
        "Row 31: Student ID O219999 not found in this batch"
      ]
    }
    
\end{lstlisting}

    Upload is processed as a background batch job. The `successCount` and
    `failCount` fields give you an immediate summary. Review any entries listed
    in `errors` and correct them before re-uploading.

    <Note>
      Uploading a file does not overwrite previously uploaded grades for students
      not included in the new file. Only rows present in the uploaded sheet are
      processed.
    </Note>

  

  
\section{Instruction Step}

    For large batches, poll the progress endpoint until the job completes.

    \begin{lstlisting}[language=bash]
    GET https://api.uniz.rguktong.in/api/v1/academics/grades/upload/progress
    Authorization: Bearer <your-token>
    
\end{lstlisting}

    **Response:**

    \begin{lstlisting}[language=json]
    {
      "success": true,
      "progress": {
        "status": "processing",
        "percent": 45,
        "processed": 150,
        "total": 330
      }
    }
    
\end{lstlisting}

    Poll every few seconds until `status` is `"completed"` and `percent` reaches
    `100`. You can also check this from the **Add Grades** page in the portal.

  
</Steps>

---

\subsection{Add grades manually for individual students}

To record or correct a grade for a single student without re-uploading a full
sheet, use the add grades endpoint directly.

\begin{lstlisting}[language=bash]
POST https://api.uniz.rguktong.in/api/v1/academics/grades/add
Authorization: Bearer <your-token>
Content-Type: application/json

\end{lstlisting}

**Request body:**

\begin{lstlisting}[language=json]
{
  "studentId": "O210008",
  "subjectCode": "E3-SEM-1-CSE-1",
  "semesterId": "SEM-1",
  "grade": "A"
}

\end{lstlisting}

**Success response:**

\begin{lstlisting}[language=json]
{
  "success": true,
  "message": "Grade added successfully"
}

\end{lstlisting}

Use this for late submissions, re-evaluations, or corrections after the main
upload has been processed.

---

\subsection{Bulk update grades via JSON}

When you need to correct multiple grades without uploading a new Excel file —
for example, after a re-evaluation — send a JSON array of updates to the bulk
update endpoint.

\begin{lstlisting}[language=bash]
PUT https://api.uniz.rguktong.in/api/v1/academics/grades/bulk-update
Authorization: Bearer <your-token>
Content-Type: application/json

\end{lstlisting}

**Request body:**

\begin{lstlisting}[language=json]
{
  "updates": [
    {
      "studentId": "O210008",
      "subjectId": "uuid-of-subject",
      "semesterId": "SEM-1",
      "grade": "EX"
    },
    {
      "studentId": "O210045",
      "subjectId": "uuid-of-subject",
      "semesterId": "SEM-1",
      "grade": 9.5
    }
  ]
}

\end{lstlisting}

The `grade` field accepts either a letter grade string (`"EX"`, `"A"`, `"F"`,
etc.) or a numeric points value like `9.5`.

<Tip>
  Use the batch grades endpoint (`GET /academics/grades/batch`) to retrieve
  subject IDs before performing bulk updates. Query params: `branch`, `year`,
  `semesterId`, and optionally `failedOnly=true` to focus on failed records.
</Tip>

---

\subsection{Download a grade report}

To download a consolidated grade report for a semester, use the download
endpoint with the semester ID.

\begin{lstlisting}[language=bash]
GET https://api.uniz.rguktong.in/api/v1/academics/grades/download/:semesterId
Authorization: Bearer <your-token>

\end{lstlisting}

**Example:**

\begin{lstlisting}[language=bash]
GET https://api.uniz.rguktong.in/api/v1/academics/grades/download/SEM-1
Authorization: Bearer <your-token>

\end{lstlisting}

The response is a downloadable report file for that semester. Use this to
verify grades before requesting publication, or to keep a local record.

---

\subsection{Publishing results to students}

<Warning>
  **Only the Director can publish results.** After you upload grades, students
  cannot see them until the Director triggers the publish action. You do not
  need to take any further steps — notify your department's administration that
  grades are ready for review and publication.
</Warning>

When the Director publishes results for a semester, UniZ triggers a background
job that sends report card emails to all students for that semester.

You can monitor the publish progress at:

\begin{lstlisting}[language=bash]
GET https://api.uniz.rguktong.in/api/v1/academics/grades/publish/progress
Authorization: Bearer <your-token>

\end{lstlisting}

**Response:**

\begin{lstlisting}[language=json]
{
  "success": true,
  "progress": {
    "status": "processing",
    "percent": 72,
    "processed": 240,
    "total": 330
  }
}

\end{lstlisting}

---

\subsection{Grade format reference}


\begin{tcolorbox}[colback=gray!5,colframe=primary,title=Documentation Note]

  
\begin{tcolorbox}[colback=gray!5,colframe=primary,title=Documentation Note]

    `EX` — Excellent

    `A` — Grade A

    `B` — Grade B

    `C` — Grade C

    `D` — Grade D

    `F` — Fail

  
\end{tcolorbox}

  
\begin{tcolorbox}[colback=gray!5,colframe=primary,title=Documentation Note]

    You can also submit grades as numeric point values (e.g., `9.5`, `7.0`).
    The system accepts both letter and numeric formats in the template and via
    the API.
  
\end{tcolorbox}

</CardGroup>


\section{Deep Technical Analysis}
The underlying system architecture for this module is built upon the \textbf{UniZ Microservices Framework}. 
This design ensures that each functional unit (Auth, Profile, Academics) remains isolated, preventing cascading failures across the university ecosystem.

\subsection{Operational Hardening}
Deployments are managed via K3s (Kubernetes) and containerized using high-performance Alpine Linux Docker images. 
Resource management is critical; for instance, the documentation service itself is provisioned with 2GB of RAM to ensure the responsive rendering of complex documentation graphs and state-management components.

\subsection{Enterprise Security Topology}
Every request entering the UniZ gateway is subjected to an aggressive JWT validation layer. 
Permissions are granular, mapping directly onto the RBAC (Role-Based Access Control) matrix defined in the core system specifications.

\subsection{Data Governance}
Prisma-driven data access ensures type safety and predictable query performance. 
We utilize PostgreSQL connection pooling to handle concurrent requests from the student portal and administrative dashboards simultaneously.
